\documentclass[10pt,a4paper]{amsart}
\usepackage[margin=19mm]{geometry}
\usepackage{amsmath,amssymb,mathtools}
\usepackage[scr=rsfs,cal=euler]{mathalfa}
\usepackage{xfrac}
\usepackage[unicode,hidelinks,bookmarks=false]{hyperref}
\newcommand{\ie}{i.e.\ }
\newcommand{\cf}{cf.\ }
\newcommand{\resp}{respectively\ }
\newcommand{\Subsection}{\subsection}
\providecommand{\nobreakdash}{\nobreak\hbox{-}}
\setlength{\emergencystretch}{2em}
\multlinegap0pt
\usepackage{mathtools}
\usepackage[all]{xy}
\usepackage{amssymb}
\usepackage{bm}
\usepackage[shortlabels]{enumitem}
\usepackage{xcolor}
\usepackage{tikz}
\newtheorem{prop}{Proposition}
\newcommand{\CH}{{\mathrm{CH}}}
\newcommand{\CHW}{{\widetilde{\mathrm{CH}}}}
\newcommand{\Er}{\mathrm{E}}
\newcommand{\K}{{{\mathbf K}}}
\newcommand{\KM}{\K^{\mathrm M}}
\newcommand{\KMW}{\K^{\mathrm{MW}}}
\newcommand{\Spec}{\operatorname{Spec}}
\title{{\bf Suslin's cancellation conjecture on real affine varieties with few real points}}
\author{Sourjya Banerjee, Jean Fasel, Samuel Lerbet}
\date{Source: arXiv:2608.17890v3 (CC BY 4.0)}
\begin{document}
\maketitle

\noindent\emph{Source and licence.} {\bf Suslin's cancellation conjecture on real affine varieties with few real points}. Version arXiv:2608.17890v3 (CC BY 4.0), distributed by arXiv under the Creative Commons Attribution 4.0 International licence, http://creativecommons.org/licenses/by/4.0/, as recorded in the arXiv licence metadata for this exact version and shown on the abstract page https://arxiv.org/abs/2608.17890v3. Full licence notice: \texttt{licenses/arxiv-2608.17890-CC-BY-4.0.txt}.\par\smallskip

\noindent\textit{Editorial scope.} Counted unit: the article's prop prop (source0.tex lines 1343--1353). The article's complete original proof of it is delivered with it (source0.tex lines 1356--1362, one \texttt{proof} environment of 7 lines). No mathematics is written, repaired or re-derived: the statement and the proof are the article's own lines, in the article's own notation, and no retained line is substituted or re-flowed.  No further source-local context is needed: every cross-reference of every retained line resolves inside the two delivered spans.  Nothing was omitted from any retained span, so the package carries no omission record.  No retained line contains a citation, so the wrapper carries no bibliography and no bibliography entry.  No retained line contains a figure environment, an image-inclusion command or drawing code, so no image is delivered and the package declares no asset dependency.  Editorial formatting only, in addition: an AMS \texttt{amsart} preamble that restates the article's own declarations and macros used by the retained lines, namely the environments \texttt{prop} on separate counters, together with the standard packages those lines need.  The retained environments print in this document as Proposition 1 in order of appearance, whereas the article prints the same items with the numbering of its own full document.  The complete native source of the article as distributed in the arXiv e-print for this exact version is bundled unchanged as \texttt{sources/207945/source0.tex}.
\begin{prop}
  Let $X=\Spec(R)$ be a smooth affine $d$-fold over a field $k$ and $2n\ge d+3$. Then, the composite
\[
\Er^n(R)\to \CHW^{n}(X)\to  \CH^{n}(X)
\]
factors through the map $\gamma^n\colon \Er^n(R)\to \Er^n_0(R)$. Consequently, there is a commutative diagram
\[
\xymatrix{\Er^{n}(R)\ar[r]\ar[d]_-{\gamma^n} & \CHW^{n}(X)\ar[d] \\
\Er_0^{n}(R)\ar[r]& \CH^{n}(X).}
\]
\end{prop}

\begin{proof}
As explained above, the composite $\Er^n(R)\to \CHW^n(X)\to \CH^n(X)$ is induced by the composite
\[
Q_{2n}\xrightarrow{\kappa} \mathrm{K}(\KMW_n,n)\to \mathrm{K}(\KM_n,n)
\]
corresponding to the class of $Z_n$ in $\CH^n(Q_{2n})$. Any element $(I,\omega_I)$ in $\Er^n(R)$ corresponds to a homotopy class of map $f\colon X\to Q_{2n}$, and its image in $\CH^n(X)$ is just $f^*(Z_n)$. The result now follows from the observation that $f^*(Z_n)$ depends only on $I$ in the pair $(I,\omega_I)$.
\end{proof}

\end{document}
